\documentclass[11pt,a4paper]{article}
\usepackage[utf8]{inputenc}
\usepackage[T1]{fontenc}
\usepackage{lmodern}
\usepackage[margin=0.95in]{geometry}
\usepackage{amsmath,amssymb,amsthm,mathtools}
\usepackage{booktabs,graphicx,array,microtype,xcolor}
\usepackage{hyperref}
\hypersetup{colorlinks=true,linkcolor=blue!60!black,citecolor=blue!60!black,urlcolor=blue!60!black}
\setlength{\emergencystretch}{3em}
\makeatletter
\DeclareRobustCommand{\_}{\ifmmode\nfss@text{\textunderscore}\else\textunderscore\allowbreak\fi}
\makeatother
\newtheorem{theorem}{Theorem}[section]
\newtheorem{proposition}[theorem]{Proposition}
\newtheorem{lemma}[theorem]{Lemma}
\newtheorem{corollary}[theorem]{Corollary}
\theoremstyle{definition}
\newtheorem{definition}[theorem]{Definition}
\newtheorem{remark}[theorem]{Remark}

\title{Exact balanced minimum evolution evaluation\\and parallel search in Phylo2Vec coordinates}
% Author names and affiliations to be supplied.
\author{}
\date{}

\begin{document}
\maketitle

\begin{abstract}
Balanced minimum evolution (BME) provides a distance-based criterion for phylogenetic inference, but searching large tree spaces requires repeated evaluation of many candidate rearrangements.
We develop a shared evaluator that computes all valid single-coordinate BME changes for an $N$-leaf tree in an ordered Phylo2Vec representation in $O(N^2)$ work and supports parallel GPU execution.
For larger trees, we optimize the connections between fixed components on a smaller skeleton tree.
Exact balanced averages of the original distances make the skeleton objective equal to the full-tree objective up to a constant.
Every skeleton change therefore gives the corresponding full-tree BME change, with an $O(K^2)$ coordinate sweep once the distances between its $K$ components are available.
Independent CPU checks agree with direct BME differences to maximum absolute discrepancies of $3.553\times10^{-14}$ for coordinate evaluation and $1.564\times10^{-13}$ for condensation.
In 120 small-tree comparisons, the CPU optimizer matches or improves the reference BME score in 119 cases.
A full coordinate sweep at $N=40{,}000$ takes 158.36~ms on eight H100 GPUs.
Exactness refers to preservation of the specified BME objective, with numerical rounding in its evaluation and heuristic selection of the search trajectory.
\end{abstract}

\section{Introduction}
\label{sec:introduction}

Distance-based phylogenetic inference first estimates pairwise evolutionary distances and then searches for a tree that explains them.
Balanced minimum evolution (BME) scores a topology using a weighted sum of those distances, with weights determined by the branching structure along each leaf-to-leaf path~\cite{pauplin2000direct,desper2002fast,desper2004theory}.
The computational challenge is to evaluate many possible rearrangements repeatedly as the tree changes.
Classical subtree-prune-and-regraft (SPR) recurrences already permit quadratic-time neighborhood sweeps~\cite{hordijk2005spr,lefort2015fastme}.
At large scales, their practical cost depends on how calculations are shared, how distance tables are stored, and how work is distributed across processors.

For a tree with $N$ leaves, we organize BME search around a table of coordinate changes in the Phylo2Vec representation.
Each entry gives the change in the objective when one integer in the tree's encoding is replaced.
The table is evaluated through shared subtree-distance calculations, exposing independent work for parallel execution.
Relabeling the current tree changes which rearrangements are represented by single-coordinate updates while preserving the tree and its objective.

The same evaluator operates on a smaller tree when some internal structure is held fixed.
We replace each fixed component by one skeleton leaf and retain its contribution through an exact balanced average of the original distances.
The full objective then equals the skeleton objective plus a constant that depends only on the fixed components.
Consequently, every change evaluated on the skeleton has the same BME change in the full tree.
We call this block of component connections a \emph{minibatch}: it restricts the structures being optimized while retaining all leaves in the objective calculation.

\begin{samepage}
We develop three algorithmic contributions:
\begin{enumerate}
 \item \textbf{Shared coordinate evaluation.} All $(N-1)^2$ valid entries are computed in $O(N^2)$ work, matching the size of the explicit table. A map from coordinate values to regraft edges avoids separately scoring candidate trees.
 \item \textbf{Exact block evaluation.} The same computation applies to a $K$-component skeleton in $O(K^2)$ work after its exact distances are available. The reduction includes the entire complement of a subtree window and specifies the updates required between windows.
 \item \textbf{Parallel execution.} Independent distance columns and SPR rows support GPU distribution and selection of the best update in each row. Reconstructing rows when needed reduces the balanced-distance table from approximately $4N^2$ to $N^2$ entries while retaining quadratic sweep work.
\end{enumerate}
\end{samepage}

The contribution builds on established representation and BME results.
Phylo2Vec introduced coordinate differences and ordered relabeling~\cite[Section~1.3.2]{penn2024phylo2vec}, and its ordered coordinate neighborhood is a subset of the SPR neighborhood.
Our evaluator retains the quadratic sweep complexity of classical BME SPR calculations.
Balanced subtree averages~\cite{desper2002fast,desper2004theory} and the fixed-clade reduction of the BME polytope~\cite[Section~2.4]{haws2011optimality} provide the basis for exact skeleton evaluation.
We state the objective identity and prove the conditions needed by the block algorithm.

GradME uses a cubic-time recurrence and automatic differentiation to optimize an expected BME objective over ordered coordinates~\cite{penn2023gradme}.
Here evaluation is specialized to a discrete tree, including its full coordinate range.
Appendix~\ref{sec:vertex_gradient} gives the connection to GradME's vertex derivatives.
Parallel taxon-addition BME~\cite{tapo2026parallel} and learned SPR selection~\cite{camerota2026learning} address tree construction and search policy, respectively.
Our experiments assess coordinate accuracy, CPU optimization, and GPU sweep throughput.

\section{Methods}
\label{sec:math}

\subsection{Overview}

The method takes a tree and a fixed matrix of distances between its leaves.
A full-tree sweep computes the objective change for every valid coordinate replacement of the current encoding.
A minibatch sweep applies the same evaluator to a skeleton of fixed components (Figure~\ref{fig:block_overview}).
After selecting and accepting a skeleton update, the components are expanded back into the full tree before the next block is constructed.
The full-tree case is recovered by taking every component to be a single leaf.

Throughout, \emph{exact} means that the calculated expression preserves the BME objective for the supplied distance matrix.
The distances themselves may be estimates, numerical evaluation uses floating-point arithmetic, and the search for an optimal topology remains heuristic.

\begin{figure}[tb]
 \centering
 \includegraphics[width=\linewidth]{figures/exact_block_overview.pdf}
 \caption{Exact block evaluation. Each colored component $S_a$ retains its internal topology and attachment root. Collapsing the components gives a skeleton whose distances $\widetilde D$ average over every original leaf pair with the appropriate balanced weights. The same matrix scores alternative skeletons, and their objective difference equals the full-tree difference after expansion. A proper subtree window includes its entire complement as one additional fixed component.}
 \label{fig:block_overview}
\end{figure}

\subsection{BME and Phylo2Vec coordinates}
\label{sec:coordinates}

Let $T$ be an unrooted binary tree on $N\geq3$ labeled leaves, and let $D$ be a fixed, finite, symmetric matrix with zero diagonal.
Write $q_T(x,y)$ for the number of internal degree-three vertices on the path from $x$ to $y$.
The BME objective is
\begin{equation}
 L(T;D)=\sum_{0\leq x<y<N}D_{xy}\,2^{-q_T(x,y)}.
 \label{eq:pauplin}
\end{equation}
If the path has $e_T(x,y)$ edges, then $q_T(x,y)=e_T(x,y)-1$.
Thus the unordered-pair coefficient is $2^{1-e_T(x,y)}$.
An artificial degree-two root used to encode the tree is suppressed before evaluating these path lengths.
No additivity assumption on $D$ is needed for the identities below.

The Phylo2Vec domain for rooted binary trees is
\begin{equation}
 \mathcal V_N=\{(v_0,\ldots,v_{N-2}):0\leq v_i\leq2i\},
 \qquad |\mathcal V_N|=(2N-3)!!.
 \label{eq:phylo2vec_space}
\end{equation}
We use the ordered, or cherry-first, subset
$\mathcal V_N^{\mathrm{cf}}=\{v:0\leq v_i\leq i\}$ as the baseline for coordinate evaluation.
Every rooted binary tree can be represented in this subset after a suitable leaf permutation.
The relabeling construction and its proof are given in Appendix~\ref{sec:gauge}.
If $\pi(a)$ is the original leaf represented by current index $a$, the distance matrix in these coordinates is
$D_\pi[a,b]=D[\pi(a),\pi(b)]$.
Permuting both axes in this way preserves the objective.
Constructing the permutation takes $O(N)$ time from a decoded tree, while explicitly copying a dense permuted matrix takes $O(N^2)$ time.

Write $T(v)$ for the unrooted tree obtained by decoding $v$ and suppressing its artificial root.
For an ordered baseline $v$, define the discrete coordinate difference
\begin{equation}
 \Delta_{i,j}
 =L\bigl(T(v^{(i\leftarrow j)});D_\pi\bigr)
  -L\bigl(T(v);D_\pi\bigr),\qquad 0\leq j\leq2i.
 \label{eq:discrete_gradient}
\end{equation}
Here $v^{(i\leftarrow j)}$ replaces only entry $i$ by $j$.
Row $i$ of the table therefore enumerates the allowed replacements of one coordinate, and $\Delta_{i,j}<0$ identifies an improving move.
We use \emph{discrete gradient} as a synonym for this table of finite objective differences, and \emph{coordinate sweep} for one evaluation of the table at a fixed tree.
There are $(N-1)^2$ valid coordinate entries, including unchanged coordinates and entries that can represent the same unrooted topology.
This count is not a count of distinct SPR neighbors.
General candidate values $j>i$ are allowed even though the baseline is ordered.
After accepting such a move, the resulting tree is decoded and relabeled before the next ordered evaluation.
Relabeling changes computational indices and permutes both axes of $D$ together, so the observations remain attached to their original leaves.

\subsection{Quadratic-time coordinate evaluation}

For two distinct vertices $u,w$, define $A(u,w)$ as the balanced average distance between the two components oriented away from the path joining $u$ and $w$.
At an internal vertex, the two outward branches each receive weight $1/2$.
At a leaf, the corresponding distribution is concentrated on that leaf.
Consequently, when the component at $u$ has children $a,b$,
\begin{equation}
 A(u,w)=\tfrac12\{A(a,w)+A(b,w)\}.
 \label{eq:bottom_up}
\end{equation}
The direction is important when one vertex is an ancestor of the other in a computational rooting.
In that case, the component away from $w$ combines the sibling branch and the parent branch:
\begin{equation}
 A(u,w)=\tfrac12\{A(\operatorname{sib}_u(w),w)
                         +A(\operatorname{parent}(u),w)\}.
 \label{eq:top_down}
\end{equation}
At the degree-three computational root, the analogous average uses the two branches not containing $w$.
Bottom-up and top-down passes construct the required table in $O(N^2)$ arithmetic operations.

For a fixed pruned component, the standard BME SPR recurrence propagates score differences across regraft edges in $O(N)$ time~\cite{hordijk2005spr}.
Each propagation step is an adjacent interchange of four directed components.
Its score change is a linear combination of their balanced distances, with coefficient $1/4$.
Successive steps accumulate the change relative to the original attachment.
Phylo2Vec supplies the mapping from each valid coordinate value to one of these regraft edges.
The explicit mapping, including the collapsed-root boundary, is in Appendix~\ref{sec:mapping}.

\begin{proposition}[Cost and scope of a coordinate sweep]
\label{prop:quadratic}
Given an ordered baseline and its leaf permutation, all entries in
Equation~\eqref{eq:discrete_gradient} can be computed in $O(N^2)$ time.
In exact arithmetic, these entries equal the corresponding full-tree BME differences.
\end{proposition}
\begin{proof}
The balanced-distance table costs $O(N^2)$.
For each of $O(N)$ coordinates, pruning its associated component and traversing the two remaining half-trees visits $O(N)$ edges.
The SPR recurrence evaluates the score at each edge, and the coordinate mapping selects the appropriate values without evaluating the candidate trees independently.
The mapping costs at most $O(N)$ per row.
The recurrence preserves the Pauplin score difference at each traversal step, so selecting a mapped edge gives Equation~\eqref{eq:discrete_gradient}.
\end{proof}

The proposition concerns the coordinate neighborhood of the current labeling.
It does not assert that this neighborhood is the entire SPR neighborhood, or that coordinate descent finds a global BME optimum.
There are $\sum_{i=0}^{N-2}(2i+1)=(N-1)^2$ valid entries.
Any algorithm that explicitly writes one scalar per entry requires $\Omega(N^2)$ work in the standard unit-cost model, so the evaluator matches this output-size lower bound.
This observation does not give a lower bound for returning only the best move, and does not improve the asymptotic sweep cost of classical BME SPR evaluation.

\subsection{Exact condensed distances and skeleton reduction}
\label{sec:condensation}

Partition all $N$ leaves into $K\geq3$ fixed components $S_1,\ldots,S_K$.
Each component has a fixed internal binary topology and a single attachment to the rest of the tree.
A component can consist of one leaf.
For a non-singleton component, let $r_a$ be the internal vertex adjacent to its attachment edge, and let $h_a(x)$ count the branching vertices from $r_a$ to leaf $x$, including $r_a$.
For a singleton, set $h_a(x)=0$.
Define its balanced leaf weights by
\begin{equation}
 \omega_a(x)=2^{-h_a(x)},\qquad x\in S_a.
 \label{eq:profile}
\end{equation}

The weights sum to one within each component (Lemma~\ref{lem:normalization} in Appendix~\ref{sec:block_proof}).

Let $G$ be an unrooted binary skeleton with one leaf for each component.
Construct $T_G$ by replacing skeleton leaf $a$ with $S_a$, retaining the attachment edge.
Define a condensed distance matrix by
\begin{equation}
 \widetilde D_{ab}
 =\sum_{x\in S_a}\sum_{y\in S_b}
    \omega_a(x)\omega_b(y)D_{xy}\quad(a\ne b),
 \qquad \widetilde D_{aa}=0.
 \label{eq:condensed_D}
\end{equation}
These are exact balanced averages of the original pairwise distances.
They can be formed by the weighted sum above or, equivalently, by recursively averaging over the two children of each component root.
All leaf contributions are retained.
This is an objective-level form of the fixed-clade reduction discussed by Haws et al.~\cite[Section~2.4]{haws2011optimality}.
We include the proof to specify the weights and the scope of exactness needed by the algorithm.

\begin{theorem}[BME decomposition over fixed components]
\label{thm:clade_condense}
There is a constant $C$, independent of the skeleton topology, such that
\begin{equation}
 L(T_G;D)=C+L(G;\widetilde D).
 \label{eq:bme_condense_decomp}
\end{equation}
Consequently, for any two skeletons $G$ and $G'$,
\begin{equation}
 L(T_{G'};D)-L(T_G;D)
 =L(G';\widetilde D)-L(G;\widetilde D).
 \label{eq:condensed_delta}
\end{equation}
\end{theorem}
The proof in Appendix~\ref{sec:block_proof} separates the fixed within-component terms from the between-component terms, whose weights factor through the skeleton.

The constant $C$ contains the contributions of pairs within the same fixed component.
It cancels when comparing skeletons, so the reduction preserves every objective change even though the skeleton's absolute score omits $C$.
The result applies to every reconnection that preserves the components and to arbitrary fixed dissimilarities, including non-additive estimated distances.

\begin{corollary}[Exact minibatch coordinate differences]
\label{cor:minibatch_gradient}
Let $v$ encode an ordered skeleton $G$, keep the attached components fixed, and index $\widetilde D$ by the same leaf labeling.
Evaluating a coordinate replacement $v_i\leftarrow j$ using the exact matrix $\widetilde D$ gives
\begin{equation}
 \Delta^{\mathrm{skel}}_{i,j}
 =L\bigl(T_{G(v^{(i\leftarrow j)})};D\bigr)
  -L\bigl(T_{G(v)};D\bigr).
 \label{eq:minibatch_gradient}
\end{equation}
Once $\widetilde D$ is available, all valid skeleton coordinate differences can be evaluated in $O(K^2)$ time.
\end{corollary}
\begin{proof}
Apply Proposition~\ref{prop:quadratic} to the $K$-leaf skeleton with input $\widetilde D$.
Equation~\eqref{eq:condensed_delta} identifies each resulting skeleton difference with the corresponding full-tree difference.
\end{proof}

\begin{corollary}[A proper subtree window]
\label{cor:outside}
For a window attached to the rest of a binary tree by one edge, treat its entire complement as an additional component oriented away from that edge.
If its internal topology and exact balanced distances to the other components are retained, Equation~\eqref{eq:condensed_delta} applies to rearrangements of the window skeleton.
\end{corollary}
\begin{proof}
Cutting the attachment edge produces two components.
The endpoint in the complement has two outward branches unless the complement is a singleton.
It therefore defines precisely a binary component of the type used in the theorem.
Including it completes the partition of all leaves, so both within-window and window-to-complement contributions are included.
\end{proof}

The outside component retains distances between leaves inside and outside the window.
Deleting the complement would change the objective being optimized.
The weights entering the outside distances are determined by the current topology of the complement and its attachment edge.
They must be updated when a preceding rearrangement changes that topology.

\subsection{Exact minibatch evaluation and sequential updates}
\label{sec:acceptance}

With the fixed components and exact condensed matrix in place, candidate skeletons are evaluated against the full-tree objective through Corollary~\ref{cor:minibatch_gradient}.
The same matrix can be reused throughout a skeleton search because its component topologies and attachment roots remain fixed.
For a candidate $G'$, let $\mathcal P(G,G')=\{(a,b):a<b,\ q_G(a,b)\ne q_{G'}(a,b)\}$.
Its full-tree BME change can also be computed directly from the changed pairs:
\begin{equation}
 \delta=\sum_{(a,b)\in\mathcal P(G,G')}
   \left(2^{-q_{G'}(a,b)}-2^{-q_G(a,b)}\right)
   \widetilde D_{ab}.
 \label{eq:changed_pairs}
\end{equation}
Unchanged pairs contribute zero, and the within-component constant cancels.
Theorem~\ref{thm:clade_condense} therefore identifies $\delta$ with the full-tree BME change.
Thus both coordinate selection and the evaluation of a combined candidate can use the same exact condensed distances.

Windows are committed sequentially.
After a commit, affected ancestor summaries are updated and the next window's outside distances are constructed from the current tree.
Node-disjoint windows need not have additive BME changes because both updates can alter the coefficients of distances between their leaves.
For consecutive accepted trees $T^{(0)},\ldots,T^{(M)}$, exact arithmetic gives
\begin{equation}
 L(T^{(M)};D)-L(T^{(0)};D)
 =\sum_{m=1}^M\bigl[L(T^{(m)};D)-L(T^{(m-1)};D)\bigr].
 \label{eq:telescoping}
\end{equation}
If each increment is negative, so is the cumulative increment.
This argument requires that each acceptance calculation uses the current tree, rather than a shared pre-update tree.

On sufficiently deep trees, finite-precision leaf weights can also underflow.

\subsection{Search strategy}
\label{sec:optimizer}

The full-tree CPU optimizer evaluates coordinate differences, ranks improving rows, and filters candidate batches for conflicts involving pruned components and attachment endpoints.
It then evaluates the combined candidate tree and backtracks to a smaller batch if the score does not improve.
The acceptance decision uses the joint objective rather than the sum of individual predicted improvements.
When the current coordinate neighborhood is exhausted, the optimizer changes the ordered labeling and searches again.
Its full-tree configuration also scans remaining unrooted SPR moves before declaring SPR-local convergence.
An optional restart stage applies random perturbations and resumes descent while retaining the best tree found.
Only accepted descent updates and the retained best score are monotone.
The perturbed working trajectory need not be.

An exact minibatch search can alternate skeletons at different scales.
A global skeleton rearranges the backbone, local windows expose detailed structure, and overlapping or paired windows propose changes between regions.
Each window follows four steps:
\begin{enumerate}
 \item Select fixed components from the current tree, including an oriented outside component for a proper subtree.
 \item Form the exact condensed distance matrix in Equation~\eqref{eq:condensed_D}.
 \item Evaluate skeleton coordinate changes using Corollary~\ref{cor:minibatch_gradient}, and check the selected update using Equation~\eqref{eq:changed_pairs}.
 \item Commit an improving candidate, graft the components back, and update summaries before selecting or evaluating the next window.
\end{enumerate}
Each accepted update decreases the same objective under the assumptions above.
Which improving moves can be reached depends on the window and relabeling schedule, so this procedure supplies no global-convergence guarantee.

\subsection{Parallel evaluation and computational cost}
\label{sec:parallel}

Equation~\eqref{eq:bottom_up} has an independent recurrence for each fixed column.
Assigning columns to threads lets neighboring threads access consecutive entries while following the same traversal order.
The top-down correction for column $w$ follows its ancestor path.
Depth-first entry intervals determine the next child without a traversal stack.
Its total work over columns is $O(\sum_w\operatorname{depth}(w))$, bounded by $O(N^2)$ and reduced to $O(N\log N)$ for trees of logarithmic height.

A full balanced-distance table stores approximately $4N^2$ entries.
The column-tiled variant retains the internal-to-leaf block and reconstructs required internal rows on demand.
This reduces that table to approximately $N^2$ entries, while the input matrix and other buffers still contribute to total memory.
A row-wise reduction can return the best coordinate replacement in each row without storing the full coordinate table.
Columns and SPR rows can also be distributed across GPUs.
Appendix~\ref{sec:parallel_details} gives the dependency and workspace arguments.

For a skeleton with $K$ leaves, a coordinate sweep costs $O(K^2)$ once its distances are available.
This does not make arbitrary distance construction independent of $N$.
Direct evaluation of one condensed pair can require $|S_a||S_b|$ original distances.
Explicit leaf weights use $O(N)$ entries, and constructing exact condensed distances can dominate runtime for large component pairs.
When $D_{xy}=\phi(x)^\top\psi(y)$ has a fixed-dimensional bilinear representation, weighted component embeddings give
$\widetilde D_{ab}=(\sum_x\omega_a(x)\phi(x))^\top
(\sum_y\omega_b(y)\psi(y))$.
That shortcut is exact for the specified bilinear distance.
It does not justify averaging embeddings before applying a nonlinear sequence-distance correction.

\section{Experimental design}
\label{sec:experiments}

The evaluation addresses three questions: whether computed changes agree with the BME objective, how the CPU optimizer behaves from a common starting tree, and how coordinate sweeps scale on GPUs.
Table~\ref{tab:experiment_scope} identifies the computation measured by each experiment.

\begin{table}[!htbp]
\centering\small
\caption{Scope of the experiments. A coordinate sweep evaluates candidates at one fixed tree. Optimization includes successive accepted updates.}
\label{tab:experiment_scope}
\begin{tabular}{@{}>{\raggedright\arraybackslash}p{.23\linewidth}>{\raggedright\arraybackslash}p{.36\linewidth}>{\raggedright\arraybackslash}p{.35\linewidth}@{}}
\toprule
Experiment & Computation & Evidence supplied\\
\midrule
Objective checks & Coordinate and condensed changes against direct BME & Numerical agreement on independently evaluated cases\\[3pt]
CPU optimization & Search from identical starting trees and distance matrices & Final objective, topology, and optimization time\\[3pt]
GPU sweeps & Full-tree coordinate evaluation & Evaluation time and memory at a fixed tree\\[3pt]
Embedding blocks & Exact bilinear condensation and skeleton sweeps & Cost of these two stages under a bilinear distance\\[3pt]
Sequence variant (Appendix~\ref{sec:sequence_variant}) & Approximate proposals followed by full-objective acceptance & Refinement time for that proposal-and-check procedure\\
\bottomrule
\end{tabular}
\end{table}

The CPU comparison uses additive distances from trees whose branch lengths are uniform on $(0.05,0.55)$ and a noisy variant
$D'_{xy}=D_{xy}\exp(0.15Z_{xy})$, with symmetric standard-normal $Z_{xy}$.
There are respectively 20, 20, 10, 5, and 5 replicates at
$N=20,50,100,150,200$ in each regime.
The reference is our in-memory C implementation of the FastME BME-SPR algorithm.
Agreement with the external FastME executable is assessed separately.
A score tolerance of $10^{-9}$ defines wins, ties, and losses.
Single-core timings are reported for Apple M5 and AMD Threadripper PRO 5995WX processors.

GPU evaluation uses NVIDIA RTX A6000 and H100 devices, with additional CPU/Metal measurements in the supplement.
The GPU coordinate path uses FP32 working arrays in the reported comparisons, whereas independent score calculations use FP64.
The embedding workload constructs condensed distances through the bilinear factorization in Section~\ref{sec:parallel}.
The sequence variant instead uses approximate condensed distances for proposals and full-tree calculations for acceptance, as detailed in Appendix~\ref{sec:sequence_variant}.
Its timings measure that procedure rather than exact condensed-distance construction for every proposal.

We assess topological error by counting inferred nontrivial splits absent from the reference tree.
For binary trees on the same leaves, normalized Robinson--Foulds (RF) distance is this count divided by $N-3$.

\section{Results}
\label{sec:results}

\subsection{Agreement with independent objective calculations}
\label{sec:validation_results}

Table~\ref{tab:validation} compares coordinate and condensed changes with independent objective calculations.
The largest discrepancy in the CPU coordinate tests is $3.553\times10^{-14}$, and the condensation tests agree to $1.564\times10^{-13}$.
These are observed absolute errors on the tested cases, not universal error bounds.
Rooted-vector round trips recover all 1,000 tested vectors at $N=40$.
Unrooted round trips preserve the split sets in a further 1,000 cases.
Exhaustive rooted-tree enumeration at $N=4,\ldots,7$ checks 11,460 encodings.

\begin{table}[tb]
\centering
\caption{Agreement with independent BME calculations. Errors are maximum absolute discrepancies over the stated tests.}
\label{tab:validation}
\small
\begin{tabular}{@{}p{0.39\linewidth}p{0.32\linewidth}r@{}}
\toprule
Calculation & Test scope & Maximum error\\
\midrule
CPU coordinate differences & 36,000 changes, $N=20$ & $3.553\times10^{-14}$\\
Ordered-tree enumeration & 28,488 changes, $N=4$--7 & $2.776\times10^{-15}$\\
Relabeling, first orientation & 200 trees, $N=25$ & $1.776\times10^{-14}$\\
Relabeling, reversed orientation & Same 200 trees & $2.665\times10^{-14}$\\
Condensed SPR differences & 16,231 changes, $N=40$, $K=16$ & $1.564\times10^{-13}$\\
FP32 CUDA gradient & Reported sizes $N=100$--4,000 & $1.11\times10^{-6}$\\
\bottomrule
\end{tabular}
\end{table}

Independent enumeration also identifies the scope of the coordinate neighborhood.
On 30 random trees at each tested size from $N=8$ to $20$, a single ordered orientation covers approximately 27--33\% of distinct SPR neighbors.
Two orientations together cover approximately 52--56\%.
Neither constitutes an exhaustive SPR search.
The full CPU optimizer's additional SPR scan is therefore material: independent enumeration found no improving SPR move in 90 returned trees at $N=10,16,24$.
This observation concerns that tested configuration, not the minibatch variant.

\subsection{CPU optimization}

Both methods recover the reference topology in all 60 noiseless trials.
On the 60 noisy trials, Phylo2Vec obtains two lower BME scores, 57 ties, and one higher score than the in-memory reference.
Thus it matches or improves the reference objective in 119 of 120 cases overall.
A better BME score need not imply a smaller RF distance when the input distances are noisy.
For example, the $N=200$ noisy comparison has a lower mean reference RF error despite one BME-score win for Phylo2Vec.

At $N=200$, the noiseless Apple M5 comparison requires 63.6 coordinate sweeps on average, versus 211.0 reference SPR iterations.
The corresponding mean times are 18.14 and 75.46~ms.
Table~\ref{tab:cpu_summary} gives the largest-size timings for both processors.
Supplementary Tables~\ref{tab:bench_noiseless} and~\ref{tab:bench_noisy} give results at all sizes.
External FastME comparisons at the same five sizes produce identical final unrooted splits and BME differences below $10^{-8}$ in the reported checks.
These checks support the reference implementation on those inputs without making the in-memory timings measurements of the external executable.

\begin{table}[tb]
\centering
\caption{Single-core optimization at $N=200$. Times are mean $\pm$ standard deviation in milliseconds over five replicates.}
\label{tab:cpu_summary}
\small
\begin{tabular}{@{}llrrr@{}}
\toprule
Processor & Distances & Reference & Phylo2Vec & Ratio\\
\midrule
Apple M5 & Additive & $75.46\pm2.42$ & $18.14\pm0.74$ & $4.16$\\
Apple M5 & Noisy & $85.16\pm2.06$ & $24.67\pm2.79$ & $3.45$\\
AMD 5995WX & Additive & $160.64\pm5.51$ & $27.96\pm1.18$ & $5.74$\\
AMD 5995WX & Noisy & $181.97\pm4.30$ & $36.91\pm3.94$ & $4.93$\\
\bottomrule
\end{tabular}
\end{table}

\subsection{Coordinate sweep throughput and memory}

At $N=4{,}000$, a full coordinate sweep takes 6.219~ms on one H100, compared with 701.375~ms for its CPU reference.
At $N=40{,}000$, distributing the full-table computation across eight H100s reduces time from 816.34 to 158.36~ms, a factor of 5.15.
The SPR stage alone falls from 756.26 to 98.28~ms.
Column tiling reduces storage at the cost of recomputation: the reported one-GPU memory requirement falls from 31,142 to 13,614~MB, while the eight-GPU tiled evaluation takes 1,032.83~ms at this size.
Thus the reduced-memory variant is useful when capacity, rather than latency, is limiting.

Distributed and streamed configurations extend a single full-tree sweep to $N=380{,}000$, taking 205.450~s with 72,946~MB reported per GPU on eight H100s.
All these times measure evaluation at a fixed tree.

For the separate embedding workload with $N=10{,}000{,}000$ and $K=20{,}000$, average condensation and coordinate-sweep times are 10.20 and 48.57~ms per batch on eight H100s.
The reported epoch estimate, 39.20~s, is the sum of those stage times multiplied by 667 batches.
It measures the projected cost of those stages over the prescribed blocks, excluding additional search work and optimization to convergence.
The complete scaling tables are in Appendix~\ref{sec:scaling_tables}.

\section{Discussion}

The method separates the size of a search block from the objective it evaluates.
A skeleton can optimize connections among a few components while retaining the contribution of every original leaf through exact condensed distances.
The same coordinate evaluator therefore supports both full-tree sweeps and restricted block searches.
This links the established BME component identity to a common computational interface for parallel tree search.

The main computational tradeoff is between reuse and reconstruction.
Dense tables permit rapid sweeps but require quadratic memory.
Column tiling trades memory for repeated arithmetic.
Skeletons reduce the dimension of the optimization problem, but exact construction of their distances can retain a dependence on the full leaf set.
This becomes particularly visible in nonlinear sequence-distance refinement and in candidates involving large condensed components.
Bilinear embedding benchmarks therefore establish a different scaling regime from the sequence experiments.

The experiments support faster coordinate evaluation and competitive CPU optimization on the tested instances.
Their scope leaves three limitations.
First, the tested relabeling and block schedules do not establish exhaustive coverage or global convergence, and the total number of sweeps remains instance-dependent.
Second, the numerical checks give observed discrepancies rather than universal error bounds, and a lower BME score under estimated distances need not imply a more accurate topology.
Third, the large sequence comparisons use approximate proposals with full-objective acceptance and do not report replicate uncertainty at the largest sizes.
Those measurements do not establish the end-to-end cost of a sequence search using exact condensed distances for every proposal.

\section{Conclusion}

Shared coordinate evaluation and exact component reduction give a common basis for full-tree and block BME search.
A quadratic sweep evaluates all valid coordinate updates, and the same computation on an exactly condensed skeleton preserves their full-tree objective changes.
The resulting parallel algorithms let a search trade the scope of its candidate set against evaluation cost while retaining the specified BME objective.

\input{phylo2vec_bme_references}

\clearpage
\appendix
\input{phylo2vec_bme_appendix}

\end{document}
